\paragrafos{Τυποποιημένη μορφή: οι μεγάλοι αριθμοί χωράνε σε μια γραμμή}{Α.3.7, «Τυποποιημένη μορφή μεγάλων αριθμών»}

\begin{anagke}
Η κυρία Αλεξάνδρα ζητάει από την Α3 μια εργασία για το ηλιακό σύστημα. Η Αγγελική γράφει στο τετράδιό της την απόσταση της Γης από τον Ήλιο: $149\,600\,000\,000$ μέτρα. Ο Ηλίας γράφει την ηλικία της Γης: $4\,500\,000\,000$ έτη, και την ξαναγράφει τρεις φορές, γιατί δεν είναι σίγουρος για τα μηδενικά. Ο Στέλιος δοκιμάζει την ηλικία του Γαλαξία και χάνεται στο δεύτερο μέτρημα. «Αυτό δεν διαβάζεται», παραπονιέται η Ρήνη.

Η Πολυξένη, που έχει το κομπιουτεράκι της, πατάει $149\,600\,000\cdot1\,000$ και η οθόνη δείχνει $1{,}496\mathrm{E}11$. «Τι είναι αυτό το E11;». Η κυρία Αλεξάνδρα χαμογελάει: «Το κομπιουτεράκι δεν χωράει όλα τα μηδενικά και τα γράφει με συντομία: $1{,}496$ επί δέκα στην ενδέκατη. Κάνουν ό,τι κάνουν και οι επιστήμονες». Η Φανή ρωτάει: «Μπορούμε να γράφουμε κι εμείς έτσι, ώστε να μη χάνουμε τα μηδενικά;».

Σε αυτή την παράγραφο θα μάθουμε να γράφουμε τους πολύ μεγάλους αριθμούς με σύντομο και ενιαίο τρόπο, τον οποίο χρησιμοποιούν οι επιστήμονες.
\end{anagke}

\begin{lexi}{τυποποιημένη μορφή}{τυποποιημένη μορφή}
Ο \lexiword{τύπος} σημαίνει καλούπι, πρότυπο. \lexiword{Τυποποιημένη} είναι η μορφή που έχει το ίδιο καλούπι για όλους τους αριθμούς: ένας αριθμός από το $1$ μέχρι κάτω από το $10$, επί μια δύναμη του $10$. Οι επιστήμονες την προτιμούν γιατί φαίνεται αμέσως το \emph{μέγεθος} του αριθμού (πόσα μηδενικά) και τα σημαντικά του ψηφία.
\end{lexi}

\begin{metafrasi}
Ένα ακολουθούμενο από $7$ μηδενικά. & $10^7=10\,000\,000$\\
Τριάντα εκατομμύρια. & $30\,000\,000=3\cdot10^7$\\
Η απόσταση Γης -- Ήλιου σε χιλιόμετρα. & $149\,600\,000=1{,}496\cdot10^8$\\
Στην οθόνη του κομπιουτεράκι. & $1{,}496\mathrm{E}8$\\
\end{metafrasi}

\tmima{Η θεωρία}

Θυμόμαστε ότι η δύναμη $10^{\nu}$ είναι ο αριθμός που γράφεται με το $1$ και μετά $\nu$ μηδενικά: $10^3=1\,000$, $10^6=1\,000\,000$.

\begin{orismos}{τυποποιημένη μορφή}
Ένας «μεγάλος» αριθμός μπορεί να γραφεί στη μορφή $\alpha\cdot10^{\nu}$, δηλαδή ως γινόμενο ενός αριθμού $\alpha$ επί μια δύναμη του $10$, όπου ο $\nu$ είναι ένας μη μηδενικός φυσικός αριθμός. Η μορφή αυτή λέγεται \textbf{τυποποιημένη}. Ο $\alpha$ είναι ένας δεκαδικός αριθμός με ακέραιο μέρος ένα ψηφίο μεγαλύτερο του $0$, δηλαδή $1\le\alpha<10$.
\end{orismos}

Η τυποποιημένη μορφή μας δίνει τη δυνατότητα να γράψουμε έναν «μεγάλο» αριθμό με απλούστερο τρόπο.

\begin{kanonas}{από τον αριθμό στην τυποποιημένη μορφή}
\begin{enumerate}[leftmargin=1.6em,itemsep=0.2ex,topsep=2pt]
\item Βάζω την υποδιαστολή μετά το \emph{πρώτο ψηφίο} (αριστερά), ώστε να προκύψει αριθμός $\alpha$ από $1$ έως $9{,}99\ldots$
\item Μετράω πόσες θέσεις μετακίνησα την υποδιαστολή προς τα αριστερά. Αυτός είναι ο εκθέτης $\nu$.
\item Γράφω $\alpha\cdot10^{\nu}$ (και αφαιρώ τα μηδενικά που περισσεύουν στο τέλος του $\alpha$).
\end{enumerate}
Για παράδειγμα $815\,000\,000\,000=8{,}15\cdot10^{11}$, γιατί η υποδιαστολή μετακινήθηκε $11$ θέσεις προς τα αριστερά. Ο εκθέτης είναι πάντα κατά $1$ μικρότερος από το πλήθος των ψηφίων του αριθμού.
\end{kanonas}

\begin{kanonas}{από την τυποποιημένη μορφή στον αριθμό}
Για να γράψω τον $\alpha\cdot10^{\nu}$ ως συνηθισμένο αριθμό, μεταφέρω την υποδιαστολή του $\alpha$ \textbf{$\nu$ θέσεις προς τα δεξιά}, συμπληρώνοντας με μηδενικά όπου χρειάζεται. Έτσι $3{,}42\cdot10^9=3\,420\,000\,000$.
\end{kanonas}

Η τυποποιημένη μορφή βοηθά και στη σύγκριση. Από δύο αριθμούς σε τυποποιημένη μορφή μεγαλύτερος είναι αυτός με τον \textbf{μεγαλύτερο εκθέτη}. Αν έχουν τον ίδιο εκθέτη, συγκρίνουμε τους $\alpha$. Έτσι $2\cdot10^{5}>9{,}9\cdot10^{4}$ και $3{,}2\cdot10^{6}<4{,}1\cdot10^{6}$.

\begin{prosoxi}
\begin{itemize}[leftmargin=1.2em,itemsep=0.3ex,topsep=2pt]
\item Το $12\cdot10^{5}$ δεν είναι τυποποιημένη μορφή, γιατί το $12$ δεν είναι μικρότερο από το $10$. Σωστά: $1{,}2\cdot10^{6}$.
\item Το $0{,}5\cdot10^{3}$ δεν είναι τυποποιημένη, γιατί το ακέραιο μέρος είναι $0$. Σωστά: $5\cdot10^{2}$.
\item Το $10^{\nu}$ δεν είναι $10\cdot\nu$. Το $10^{3}=1\,000$ και όχι $30$.
\item Ο εκθέτης δεν είναι το πλήθος των ψηφίων αλλά το πλήθος των μηδενικών μετά το $1$, δηλαδή τα ψηφία μείον $1$: ο $5\,000\,000$ έχει $7$ ψηφία και είναι $5\cdot10^{6}$.
\item Τα μηδενικά στο τέλος του $\alpha$ δεν γράφονται: $2{,}50\cdot10^4=2{,}5\cdot10^4$.
\end{itemize}
\end{prosoxi}

\tmima{Λυμένα παραδείγματα}

\paradeigma{Στην τυποποιημένη μορφή}
\tixero Οι αριθμοί $30\,000\,000$ και $815\,000\,000\,000$.\par
\tithelo Να τους γράψω σε τυποποιημένη μορφή.\par
\possindeo Ο πρώτος έχει $8$ ψηφία: μετακινώ την υποδιαστολή κατά $7$ θέσεις και παίρνω $3$. Ο δεύτερος έχει $12$ ψηφία: κατά $11$ θέσεις, και παίρνω $8{,}15$.\par
\apantisi $30\,000\,000=3\cdot10^{7}$ και $815\,000\,000\,000=8{,}15\cdot10^{11}$.\par
\elegxos $3\cdot10^{7}=3\cdot10\,000\,000$ και $8{,}15\cdot10^{11}$: η υποδιαστολή πάει $11$ θέσεις δεξιά, και παίρνουμε το $815$ και $9$ μηδενικά.

\paradeigma{Από την τυποποιημένη μορφή}
\tixero Ο αριθμός $3{,}42\cdot10^{9}$.\par
\tithelo Να τον γράψω ως συνηθισμένο αριθμό.\par
\possindeo Μετακινώ την υποδιαστολή κατά $9$ θέσεις δεξιά. Το $3{,}42$ έχει δύο δεκαδικά ψηφία, οπότε προσθέτω $7$ μηδενικά.\par
\apantisi $3\,420\,000\,000$.\par
\elegxos Ο αριθμός έχει $10$ ψηφία, ενώ ο εκθέτης είναι $9$.

\paradeigma{Σύγκριση}
\tixero Οι αριθμοί $2\cdot10^{3}$ και $3\cdot10^{2}$, και οι $0{,}1\cdot10^{2}$ και $10$.\par
\tithelo Να τους συγκρίνω.\par
\possindeo Στο πρώτο ζεύγος οι εκθέτες διαφέρουν: $2\,000>300$. Το δεύτερο ζεύγος: το $0{,}1\cdot10^{2}$ δεν είναι τυποποιημένο. Είναι $0{,}1\cdot100=10$.\par
\apantisi $2\cdot10^{3}>3\cdot10^{2}$ και $0{,}1\cdot10^{2}=10$.\par
\elegxos $2\cdot10^3=2\,000$ και $3\cdot10^2=300$.

\paradeigma{Χρόνος σε λεπτά}
\tixero Οι καρχαρίες της Γροιλανδίας ζουν περίπου $400$ χρόνια. Ένα έτος έχει περίπου $365$ ημέρες.\par
\tithelo Πόσα λεπτά είναι τα $400$ χρόνια; Να δώσω την απάντηση σε τυποποιημένη μορφή.\par
\possindeo Ένα έτος έχει $365\cdot24\cdot60=525\,600$ λεπτά. Τα $400$ έτη έχουν $400\cdot525\,600=210\,240\,000$ λεπτά. Ο αριθμός έχει $9$ ψηφία, άρα ο εκθέτης είναι $8$.\par
\apantisi $2{,}1024\cdot10^{8}$ λεπτά.\par
\elegxos $2{,}1024\cdot10^8$: η υποδιαστολή πάει $8$ θέσεις δεξιά και δίνει $210\,240\,000$.

\paradeigma{Ποιοι αριθμοί είναι τυποποιημένοι;}
\tixero Οι $4{,}5\cdot10^{3}$, $45\cdot10^{2}$, $0{,}45\cdot10^{4}$, $1\cdot10^{6}$.\par
\tithelo Να βρω ποιοι είναι σε τυποποιημένη μορφή και να διορθώσω τους άλλους.\par
\possindeo Ελέγχω αν το ακέραιο μέρος του πρώτου παράγοντα είναι ένα ψηφίο από το $1$ ως το $9$. Στο $4{,}5$ και στο $1$ ναι. Το $45$ έχει δύο ψηφία και το $0{,}45$ έχει ακέραιο μέρος $0$.\par
\apantisi Τυποποιημένοι είναι οι $4{,}5\cdot10^{3}$ και $1\cdot10^{6}$. Οι άλλοι δύο είναι ίσοι με τον πρώτο, $45\cdot10^2=4\,500=0{,}45\cdot10^4=4{,}5\cdot10^3$.\par
\elegxos Και οι τρεις δίνουν $4\,500$.

\tmima{Ασκήσεις}

\zoni{Το κατάλαβα}{από και προς την τυποποιημένη μορφή}

\ex Γράψε σε τυποποιημένη μορφή: α) $150\,000$ \quad β) $21\,000$ \quad γ) $1\,300\,000$ \quad δ) $2\,030\,000$.

\ex Γράψε σε πλήρη μορφή: α) $2{,}31\cdot10^{4}$ \quad β) $1{,}36\cdot10^{9}$ \quad γ) $1{,}37\cdot10^{2}$ \quad δ) $10^{6}$.

\ex Σύγκρινε τους αριθμούς με ένα από τα $<,\ >,\ =$: α) $2\cdot10^{3}\ \ldots\ 3\cdot10^{2}$ \quad β) $0{,}1\cdot10^{2}\ \ldots\ 10$ \quad γ) $1{,}7\cdot10^{4}\ \ldots\ 3\cdot10^{3}$ \quad δ) $6\cdot10^{5}\ \ldots\ 0{,}6\cdot10^{6}$.

\ex Ποιοι από τους αριθμούς $4{,}5\cdot10^{3}$, $45\cdot10^{2}$, $0{,}45\cdot10^{4}$, $9{,}99\cdot10^{8}$, $10\cdot10^{5}$, $1\cdot10^{6}$ είναι σε τυποποιημένη μορφή;

\ex Γράψε σε τυποποιημένη μορφή τους αριθμούς που δεν είναι, από την προηγούμενη άσκηση.

\ex Γράψε στη μορφή $\alpha\cdot10^{\nu}$ τους αριθμούς: α) $7\,000$ \quad β) $480\,000\,000$ \quad γ) $12\,500\,000\,000$ \quad δ) $6\,021\,000$.

\ex Σωστό ή λάθος; α) $3\cdot10^{5}=300\,000$. \quad β) $2{,}5\cdot10^{4}=250\,000$. \quad γ) Ο $10^{7}$ γράφεται με το $1$ και επτά μηδενικά. \quad δ) $5{,}2\cdot10^{3}>4{,}9\cdot10^{3}$. \quad ε) $3\cdot10^{2}=6$. \quad στ) $7\cdot10^{1}=70$.

\zoni{Το χρησιμοποιώ}{διάστημα, χρόνος, πληθυσμοί}

\ex \textbf{Τα αστρονομικά μεγέθη.} Γράψε σε τυποποιημένη μορφή: α) Απόσταση Γης -- Σελήνης: $384\,400$ χιλιόμετρα. \quad β) Απόσταση Γης -- Ήλιου: $149\,600\,000$ χιλιόμετρα. \quad γ) Ηλικία της Γης: $4\,500\,000\,000$ έτη. \quad δ) Ηλικία του Γαλαξία: $13\,600\,000\,000$ έτη. \quad ε) Απόσταση Πλούτωνα -- Ήλιου: $5\,900\,000\,000$ χιλιόμετρα.

\ex \textbf{Η χελώνα.} Μια γιγάντια χελώνα ζει περίπου $150$ χρόνια. Σε πόσα λεπτά αντιστοιχούν (με έτος $365$ ημερών); Γράψε την απάντηση σε τυποποιημένη μορφή.

\ex \textbf{Οι πλανήτες.} Οι αποστάσεις από τον Ήλιο: Δίας $7{,}8\cdot10^{8}$ km, Άρης $2{,}3\cdot10^{8}$ km, Ποσειδώνας $4{,}5\cdot10^{9}$ km, Ερμής $5{,}8\cdot10^{7}$ km. Γράψε τους πλανήτες από τον πιο κοντινό στον πιο μακρινό.

\ex \textbf{Το κομπιουτεράκι.} Το κομπιουτεράκι της Πολυξένης δείχνει $3{,}2\mathrm{E}8$ και μετά $4{,}05\mathrm{E}6$. Ποιοι είναι οι αριθμοί;

\ex \textbf{Οι πληθυσμοί.} Η Κίνα έχει περίπου $1\,400\,000\,000$ κατοίκους και η Ελλάδα περίπου $10\,400\,000$. Γράψε τους αριθμούς σε τυποποιημένη μορφή και πες ποιος είναι μεγαλύτερος και πώς το βλέπεις από τους εκθέτες.

\zoni{Το σκέφτομαι}{γιατί; και τι θα συνέβαινε αν...;}

\ex Εξήγησε γιατί το $12\cdot10^{5}$ δεν είναι σε τυποποιημένη μορφή και γράψε το σωστά.

\ex Πόσα ψηφία έχει ένας αριθμός που γράφεται στη μορφή $\alpha\cdot10^{\nu}$; Πόσα ψηφία έχουν οι $4\cdot10^{9}$ και $4{,}5\cdot10^{9}$;

\ex Ο Στέλιος έγραψε $5\,400\,000=5{,}4\cdot10^{5}$. Βρες το λάθος του.

\ex Ανάμεσα σε ποιες δύο διαδοχικές δυνάμεις του $10$ βρίσκεται ο $4{,}5\cdot10^{6}$;

\zoni{Για το σπίτι}{ψάξε και γράψε}

\ex Ψάξε πέντε μεγάλα μεγέθη από την εγκυκλοπαίδεια ή το διαδίκτυο (π.χ. πληθυσμός της Γης, μάζα ή ακτίνα πλανήτη) και γράψε τα σε τυποποιημένη μορφή.

\ex Υπολόγισε πόσα δευτερόλεπτα έχει ένα έτος ($365$ ημέρες) και γράψε το αποτέλεσμα σε τυποποιημένη μορφή.
